%% kapa: exam=1 round=22 year=2011 period=1 form=A subject=economics subject_ko=경제원론 seq=12 number=52
%% source: Q-Net 감정평가사 기출문제 (cst00309) artlSeq=5036790 file=extracted/1008/5036790/2011년 제22회 감정평가사 1차시험 기출문제/1교시.hwp
%% conversion: kordoc 4.24.51 -> tools/gen1.py (auto); review=proofread-sonnet; flags=none
\begin{question}{52}
경제성장모형에서 생산함수가 Y=AK 일 때 다음 설명 중 옳은 것만을 모두 고른 것은? (단, Y는 생산량, A는 생산성수준이며 0보다 큰 상수, K는 자본량)
\begin{framedtext}
ㄱ. 자본의 한계생산물은 일정하다. \par
ㄴ. 자본량이 증가할 때 생산량은 증가한다. \par
ㄷ. 노동량이 증가할 때 생산량은 증가한다. \par
ㄹ. 자본의 증가율과 생산량의 증가율은 같다.
\end{framedtext}
\begin{choices}
  \item ㄱ, ㄴ
  \item ㄱ, ㄴ, ㄹ
  \item ㄱ, ㄷ, ㄹ
  \item ㄴ, ㄷ, ㄹ
  \item ㄱ, ㄴ, ㄷ, ㄹ
\end{choices}
\end{question}
\answer{2}
